% Appendix C: Six-port incidence and orbit counts
\label{app:six-port-counts-body}

This appendix records the finite counts used in \S\ref{sec:apertures},
\S\ref{sec:six-port}, and Appendix~\ref{app:proofs}.  Where a count is stated in the source
material, the computation is reproduced here so that it can be checked independently.

\subsection{The support partition}

The carrier of Proposition~\ref{prop:seven-faces} has seven nonempty faces: three vertices,
three edges, and one interior.  Its $2^{7}=128$ supports partition by cardinality as
\[
  128=\underbrace{2}_{|S|\in\{0,7\}}+\underbrace{6}_{|S|\in\{1,6\}}
  +\underbrace{30}_{|S|\in\{2,5\}}+\underbrace{90}_{|S|\in\{3,4\}} .
\]
The two-element classes are the empty and full supports; the six-element classes are the
singletons and their complements; the remaining two classes are self-complementary in the
sense that a support of size $3$ has a complement of size $4$.

The refinement $30=18+12$ separates the supports of size $2$ and $5$ into those fixed by an
orientation-reversing reflection and those exchanged by it; the $18$ are achiral and the $12$
chiral.

\subsection{The $D_3$ orbit count}

The dihedral group $D_3$ of order $6$ acts on the seven faces: it permutes the three vertices
in a $3$-cycle or a transposition, permutes the three edges identically, and fixes the
interior.  Burnside's lemma on the induced action on the $128$ supports gives:

\begin{center}
\begin{tabular}{@{}lcl@{}}
\toprule
Element & Number & Fixed supports \\
\midrule
identity & $1$ & $128$ \\
$120^\circ$ rotation & $2$ & $2\cdot2\cdot2=8$ each \\
reflection & $3$ & $4\cdot4\cdot2=32$ each \\
\bottomrule
\end{tabular}
\end{center}

A rotation acts as a $3$-cycle on the vertices and a $3$-cycle on the edges, so a fixed support
is determined by $2$ choices per $3$-cycle and $2$ choices for the interior: $8$.  A reflection
fixes one vertex and swaps a pair, and likewise for edges, so a fixed support has
$2\cdot2$ choices on vertices and $2\cdot2$ on edges and $2$ on the interior: $32$.
Hence
\[
  \#\text{orbits}=\frac{1}{6}\bigl(128+2\cdot8+3\cdot32\bigr)
  =\frac{240}{6}=40,
\]
reproducing the count of Proposition~\ref{prop:seven-faces}.  Pairing orbits under the
negation involution $C_2^{\neg}$ halves this to the $20$ archetypes stated there.

\subsection{The matching count for the bridge}

Theorem~\ref{thm:m6-bridge} uses two fixed-point-free involutions on six elements.  The exact
counts are as follows.

\begin{proposition}[Matching counts]\label{prop:matching-counts}
On a six-element set there are
\begin{enumerate}
\item $15$ fixed-point-free involutions (perfect matchings);
\item $15\cdot14=210$ ordered pairs of distinct such involutions;
\item $60$ six-cycles, each admitting exactly $2$ alternating ordered $2$-colourings;
\item hence $120$ ordered pairs whose union is connected, and $210-120=90$ whose union is two
  disjoint triangles.
\end{enumerate}
\end{proposition}

\begin{proof}
(i) A perfect matching of six elements is a partition into three unordered pairs, of which
there are $\frac{6!}{2^{3}\,3!}=15$.  (ii) Immediate.  (iii) The number of cycles on six
labelled vertices up to rotation and reflection is $(6-1)!/2=60$, and a cycle of even length
has exactly two alternating $2$-colourings, one for each choice of starting colour.
(iv) $60\cdot2=120$; the rest are unions of two triangles, which arise when the two matchings
are the alternating classes of two disjoint $3$-cycles.
\end{proof}

The count $120$ is the exact size of the family to which
Theorem~\ref{thm:m6-bridge} applies: connected unions are the case in which the presented
boundary is a single $M_6^{\partial}$ rather than a disjoint union of smaller boundaries.

\subsection{Miscellaneous finite counts}

\begin{itemize}
\item The declared maps from six positions into three domains number $3^{6}=729$; the
  retention ledger of the grammar is indexed by these.
\item The set partitions of six positions number $B_6=203$ (the sixth Bell number), which is
  the count of restricted-growth structures used for the allocation ledger.
\item The permutohedron on four symbols has $24$ vertices, matching the $TO_{24}$ vertex
  count of Proposition~\ref{prop:envelope}; the two are different graphs with the same vertex
  count, and the paper does not identify them.
\item The relation boundary count is $2m$; for $m=3$ this is $6$, which is the port count of
  $\bW_6$ and the degree $\deg(\bW_6)$ of Definition~\ref{def:whole-cut-record}.
\end{itemize}

\subsection{What these counts do and do not establish}

Every count in this appendix is an exact statement about a declared finite fixture.  None of
them is a theorem about arbitrary carriers, and in particular:

\begin{itemize}
\item the $40$ orbits and $20$ archetypes are properties of the seven-face triangle, not of
  three-domain carriers in general;
\item the $120$ connected unions count presentations, not fillers: the filler remains open by
  Calibration~\ref{cal:m6-open};
\item the equality of the $TO_{24}$ and permutohedron vertex counts is a numerical coincidence
  unless a map is supplied, and none is supplied here.
\end{itemize}
